\section{Safety Discussion}
\label{sec:safety}

The induction circuit is fundamental to in-context learning~\citep{olsson2022context}.
A common a priori concern is that fine-tuning on a narrow distribution might
\emph{degrade} this circuit, causing deployed fine-tuned models to lose
in-context generalisation outside their target domain---a property standard
capability benchmarks do not measure.

\textbf{Implication for fine-tuning safety evaluations.}
Our preliminary result (Section~\ref{sec:results}) does not support the
degradation hypothesis at this token budget: L1H6's prefix-matching score
\emph{increased} under both code ($+0.194$ at step 100, $+0.242$ cumulative
by step 200) and prose ($+0.175$ at step 100, $+0.209$ cumulative) fine-tuning.
If this preliminary finding holds across the remaining seeds, it would
suggest that short-horizon fine-tuning on the distributions tested here is
more likely to \emph{reinforce} prefix-matching behaviour than to silently
erode it. This is itself a safety-relevant finding: a model that becomes
\emph{more} reliant on a simple copy-based heuristic after fine-tuning may be
more susceptible to prompt-injection or repetition-based attacks that exploit
that heuristic, even as it appears unchanged or improved on standard
benchmarks. Either direction of change---degradation or reinforcement---is
invisible to capability benchmarks alone, which is the core argument for
including circuit-level diagnostics in fine-tuning evaluations regardless of
which direction the change turns out to go.

\textbf{Phase transitions and early stopping.}
The detected transition occurs predominantly in the first 100 of 244 steps
in both conditions, with comparatively little further change between steps
100 and 200. If this front-loaded pattern generalises, it implies that
circuit-level monitoring is most informative early in a fine-tuning run, and
that validation-loss-based early stopping criteria---which are insensitive to
attention-pattern changes---would not by themselves indicate whether or when
such a change has occurred.
